\documentclass[/work/src/main.tex]{subfiles}

\begin{document}

\section{Contenido Principal}

\subsection{Guía de Estructura}
Este documento está organizado de forma modular utilizando el paquete \texttt{subfiles}. Cada sección se encuentra en un archivo independiente dentro de la carpeta \texttt{sections/}, permitiendo:

\begin{itemize}
    \item Compilación independiente de cada sección
    \item Mejor organización del código
    \item Facilidad de mantenimiento y edición
    \item Reutilización de paquetes y configuración
\end{itemize}

\subsection{Configuración del Proyecto}
La configuración general del documento se encuentra centralizada en la carpeta \texttt{settings/}, que incluye:

\begin{itemize}
    \item \textbf{Paquetes:} definición de todos los paquetes LaTeX necesarios
    \item \textbf{Estilos:} colores, fuentes y diseño visual
    \item \textbf{Bibliografía:} gestión de referencias bibliográficas
    \item \textbf{Macros:} comandos personalizados reutilizables
    \item \textbf{Teoremas:} definición de teoremas, lemas y proposiciones
    \item \textbf{Metadatos:} información del documento (título, autor, fecha)
\end{itemize}

\subsection{Mejoras Tipográficas}
El documento implementa varias mejoras tipográficas automáticas:

\subsubsection{microtype}
Mejora sutil pero notablemente el espaciado y la justificación del texto mediante ajustes microtipográficos. Se aplica automáticamente a todo el documento sin necesidad de intervención.

\subsubsection{cleveref}
Permite escribir referencias cruzadas inteligentes. Por ejemplo, \texttt{\textbackslash cref\{etiqueta\}} genera automáticamente ``Figura 1'' o ``Sección 2'' detectando el tipo de objeto.

\subsubsection{csquotes}
Maneja comillas de forma inteligente según el idioma especificado en Babel. Con babel en español, las comillas se adaptan automáticamente: \enquote{así se ve una comilla inteligente}.

\subsection{Sistema Bibliográfico}
El documento utiliza \textbf{BibLaTeX} con el backend \texttt{biber}, un sistema moderno de gestión bibliográfica que simplifica las referencias:

\begin{itemize}
    \item Soporta UTF-8 sin problemas
    \item Fácil de personalizar estilos
    \item Carga automática desde \texttt{bibliography.bib}
    \item Mucho más flexible que BibTeX clásico
\end{itemize}

Para citar una referencia, simplemente usa \texttt{\textbackslash cite\{clave\}}. Por ejemplo, \cite{greenwade93} es una referencia a un artículo sobre CTAN. La bibliografía completa aparecerá automáticamente al final del documento.

\end{document}
